\documentclass[11pt]{article}
\usepackage{amsmath,amssymb,booktabs,geometry}
\geometry{margin=1in}
\newcommand{\Tr}{\operatorname{Tr}}
\newcommand{\Sym}{\operatorname{Sym}}
\newcommand{\Ran}{\operatorname{Ran}}
\begin{document}
\title{Nine symmetric qubits: the maximum second-order anticoherent subspace dimension is one}
\author{Exact catalogue calculation}
\date{}
\maketitle
\paragraph{Theorem.}
For pure symmetric nine-qubit subspaces,
\[
 \boxed{k_{\max}(9,2)=1.}
\]
The proof applies to arbitrary complex subspaces. Its finite exact certificate is stored in \texttt{kernel.json} and \texttt{certificate.json}; the accompanying verifier reconstructs the representation and checks the rational identities.

\section*{1. The unrestricted complex problem}
Let $V$ be the spin-$9/2$ representation, of dimension ten. Suppose $P$ projects onto a two-dimensional pure symmetric second-order anticoherent subspace. If $u,v$ are an arbitrary complex orthonormal basis of $\Ran P$, then
\[
 PT_{Lq}P=0\qquad (L=1,2;\ -L\le q\le L),
\]
where $T_{Lq}$ are irreducible spin tensors. This is the full compressed-operator condition, including all cross terms between $u$ and $v$.

Define the normalized exterior product
\[
 w=u\wedge v=\frac{u\otimes v-v\otimes u}{\sqrt2},
 \qquad R_w=|w\rangle\langle w|.
\]
For any operator $T$ on $V$, direct expansion gives
\begin{align*}
2\Tr_1[(T\otimes I)R_w]
={}&\langle u|T|u\rangle|v\rangle\langle v|
+\langle v|T|v\rangle|u\rangle\langle u|\\
&-\langle v|T|u\rangle|v\rangle\langle u|
-\langle u|T|v\rangle|u\rangle\langle v|.
\end{align*}
Since $u,v$ are orthonormal, this matrix vanishes if and only if $PTP=0$. In particular,
\[
 \Tr_1[(T_{Lq}\otimes I)R_w]=0\quad(L=1,2).
\]

\section*{2. Four copies and the Young subspace}
Set $\omega=w\otimes w\in V^{\otimes4}$. It lies in
$\Sym^2(\Lambda^2V)$, and its component in $\Lambda^4V$ vanishes because
$w\wedge w=0$. Therefore
\[
 \omega\in W:=\mathbb S_{(2,2)}V,
 \qquad \dim W=\frac{10^2(10^2-1)}{12}=825.
\]
An explicit orthogonal projector onto $W$ is
\[
 Y_{22}=E-A_4,\qquad
 E=\frac{I-U_{(12)}}2\frac{I-U_{(34)}}2
       \frac{I+U_{(13)(24)}}2,
 \qquad A_4=\frac1{24}\sum_{\pi\in S_4}\operatorname{sgn}(\pi)U_\pi.
\]
Here $A_4$ is a subprojector of $E$, so $Y_{22}$ is indeed an orthogonal projector. This construction depends only on permutation symmetries, not on an ansatz for the code.

Let $D(g)$ denote the spin-$9/2$ representation and define the collective rotation average
\[
 X=\int_{SU(2)}D(g)^{\otimes4}|\omega\rangle\langle\omega|
                      D(g)^{\dagger\otimes4}\,dg.
\]
Then $X\succeq0$, $\Tr X=1$, $\Ran X\subseteq W$, and $X$ commutes with collective rotations. Moreover,
\[
 \mathcal C_{Lq}(X):=
 \Tr_1[(T_{Lq}\otimes I^{\otimes3})X]=0,
 \qquad L=1,2.
\]
Before averaging, this identity is the tensor product of the zero matrix in Section~1 with $R_w$. Averaging preserves the identity because the tensors of a fixed rank form a rotation-invariant linear space.

For a rotation-invariant $X$, it suffices to impose $q=0$. Indeed,
$T\mapsto\Tr_1[(T\otimes I^{\otimes3})X]$ intertwines rotations, and the orbit of the nonzero vector $T_{L0}$ spans the entire irreducible rank-$L$ tensor space. Thus using only $T_{10}$ and $T_{20}$ in the linear system loses no rank-one or rank-two contraction condition.

\section*{3. Completeness of the finite representation}
Under collective rotations,
\[
 W\cong\bigoplus_J V_J\otimes\mathbb C^{m_J}.
\]
The multiplicities are
\[
\begin{array}{c|rrrrrrrrrrrrrrr}
J&0&2&3&4&5&6&7&8&9&10&11&12&13&14&16\\\hline
m_J&3&6&2&7&4&7&4&7&3&5&2&3&1&2&1.
\end{array}
\]
They satisfy
\[
 \sum_J(2J+1)m_J=825,
 \qquad \sum_J\frac{m_J(m_J+1)}2=169.
\]
The table can be obtained independently from
$\Sym^2(\Lambda^2V)-\Lambda^4V$: count each weight, then subtract consecutive nonnegative weight multiplicities to recover $m_J$. An exact highest-weight construction must find exactly $m_J$ independent vectors in
$W\cap\ker J_+\cap\ker(J_z-JI)$ for every listed $J$. Their lowering orbits then span all 825 dimensions. These dimension checks prevent accidental omission of a possible subspace.

In an orthonormal real coupled basis, Schur's lemma gives
\[
 X=\bigoplus_J\frac{I_{2J+1}}{2J+1}\otimes\sigma_J,
 \qquad \sigma_J\succeq0,\quad\sum_J\Tr\sigma_J=1.
\]
The matrices $\sigma_J$ may initially be complex Hermitian. The tensor matrices $T_{L0}$, the permutation projector, and the chosen coupled basis have real entries. Consequently the contraction equations have real coefficients. Replacing every $\sigma_J$ by
\[
 \frac{\sigma_J+\overline{\sigma_J}}2
\]
preserves positivity, trace, and all contraction equations. Thus a hypothetical complex code always supplies real symmetric positive semidefinite blocks. This reduction does \emph{not} assume that the original states have real Dicke coefficients.

\section*{4. The exact coordinates and their normalization}
To avoid unnecessary square roots, the saved calculation uses real, unnormalized highest-weight columns $H_{Ja}$, $a=1,\ldots,m_J$. They need not be orthogonal. Put
\[
 H_{JaM}=\frac{J_-^{J-M}}{(J-M)!}H_{Ja},\qquad
 (G_J)_{ab}=\langle H_{Ja},H_{Jb}\rangle.
\]
The ladder identity gives
\[
 \langle H_{JaM},H_{JbM}\rangle
 =\binom{2J}{J-M}(G_J)_{ab}.
\]
Every invariant operator on $W$ can therefore be written uniquely as
\[
 X=\sum_{J,M}\sum_{a,b=1}^{m_J}
 C^J_{ab}\,
 \frac{|H_{JaM}\rangle\langle H_{JbM}|}{\binom{2J}{J-M}}.
\]
The coefficient matrices $C^J$ are real symmetric. Because each highest-weight matrix has full column rank, changing between $C^J$ and orthonormal multiplicity blocks is an invertible congruence. Consequently
\[
 X\succeq0\quad\Longleftrightarrow\quad C^J\succeq0
 \ \hbox{for every }J.
\]
In particular, all 57 diagonal entries $C^J_{aa}$ are nonnegative. The trace is
\[
 \Tr X=\sum_J(2J+1)\Tr(G_JC^J).
\]
It generally contains off-diagonal terms because $G_J$ need not be diagonal. The proof below does not replace this trace by an unweighted sum of block diagonals.

Use the 169 upper-triangular entries of the $C^J$ as coordinates $x$. An off-diagonal coordinate contributes both symmetric matrix entries. Exact expansion of selected entries of the two contraction matrices gives
\[
 Ax=0,\qquad A\in\mathbb Q^{156\times169},\qquad\operatorname{rank}A=156.
\]
Only a subset of the full contraction equations is required; every genuine code satisfies this subset. The stored exact kernel basis
\[
 K\in\mathbb Q^{169\times13},\qquad AK=0,
 \qquad\operatorname{rank}K=13
\]
therefore spans all solutions of these necessary equations. Let $D$ select the 57 diagonal coordinates from $x$.

\section*{5. A strictly positive exact certificate}
The file \texttt{certificate.json} gives a vector $z\in\mathbb Z^{57}$, with every entry strictly positive, such that
\[
 \boxed{z^{\mathsf T}DK=0.}
\]
Its largest coefficient has 106 decimal digits. No rounded coefficient or numerical sign decision enters this identity. The certificate is deliberately supplied as machine-readable integers rather than several pages of large numerators.

For any solution $x=Ky$ of the necessary contraction equations,
\[
 \sum_{J,a}z_{Ja}C^J_{aa}=z^{\mathsf T}Dx=0.
\]
Since every $z_{Ja}>0$ and $C^J\succeq0$, all diagonal entries vanish. For a positive semidefinite matrix,
\[
 |C^J_{ab}|^2\le C^J_{aa}C^J_{bb},
\]
so every off-diagonal entry also vanishes. Thus every $C^J=0$, and hence $X=0$. This contradicts $\Tr X=1$.

The exact verification includes all 2028 entries of $AK$, all 13 entries of $z^{\mathsf T}DK$, all 57 strict signs, and the representation dimensions. The ranks can be independently certified by nonzero minors modulo the prime $1000003$; all rational denominators are nonzero modulo that prime. The physical row reconstruction separately verifies that the saved rows are the stated contraction equations on the complete Young subspace.

This excludes every two-dimensional second-order anticoherent subspace. A higher-dimensional subspace would contain a two-dimensional one, so it is excluded as well.

\section*{6. Exact lower bound}
Dimension one is attained, for example by
\[
 |\psi\rangle=\frac{\sqrt{13}|D_9^{(0)}\rangle
                         +\sqrt{30}|D_9^{(4)}\rangle
                         -\sqrt{21}|D_9^{(8)}\rangle}{8}.
\]
It is normalized, its Dicke support gaps eliminate all nonzero magnetic tensor components through rank two, and
$\langle J_z\rangle=0$, $\langle J_z^2\rangle=33/4=j(j+1)/3$.
Its span is therefore second-order anticoherent. Together with the exclusion in Section~5, this establishes $k_{\max}(9,2)=1$.
\end{document}
